% !TEX root = ../main.tex
\chapter{Supplementary Comparison Figures}
\label{app:supplementary_comparison_figures}

This appendix collects supplementary figures for the spectroscopy analysis in
the main chapter. It includes monomer validation and waiting-time figures, as
well as dimer control panels and waiting-time series used alongside the main
discussion.

Unless stated otherwise, the supplementary signal figures follow the layout
convention introduced at the beginning of~\autoref{chapter:results}. Axis
labels and, where present, colour-bar labels apply to all axes and colour bars
in the corresponding figure.

\section{Monomer Figures}
\label{app:sec:paper_eqs_monomer}

\subsection{Validation Figures}

This subsection collects the monomer validation figures used throughout the
thesis. It includes the inhomogeneous convergence sweeps and the higher-field
sensitivity checks.

\begin{figure}[htbp]
    \centering
    \safeplaceholdergraphic[1.22\linewidth]{figures/app_paper_eqs_monomer_fig2_time}
    \caption{Monomer 1D inhomogeneous time-domain convergence figures for
    $N_{\mathrm{inhom}}=10$ (a), $100$ (b), $1000$ (c), and $10000$ (d).
    Shared settings are listed in~\autoref{subsec:results_monomer_model}.}
    \label{fig:app_paper_eqs_monomer_fig2_time_n_dep}
\end{figure}

\begin{figure}[htbp]
    \centering
    \safeplaceholdergraphic[1.22\linewidth]{figures/app_paper_eqs_monomer_fig2_freq}
    \caption{Monomer 1D inhomogeneous frequency-domain convergence figures for
    the same $N_{\mathrm{inhom}}$ sweep as
    in~\autoref{fig:app_paper_eqs_monomer_fig2_time_n_dep}. Shared settings
    are listed in~\autoref{subsec:results_monomer_model}.}
    \label{fig:app_paper_eqs_monomer_fig2_freq_n_dep}
\end{figure}


\subsection{Waiting-Time Figures}

The monomer waiting-time behaviour is discussed in the main chapter.
\autoref{fig:rslts_monomer_2d_time_hom_inhom} shows the \(t_{\mathrm{wait}}=0\)
reference case in the 2D time domain, and
\autoref{fig:rslts_monomer_2d_freq_hom_inhom_rephasing} shows the corresponding
rephasing waiting-time series in the 2D frequency domain. This appendix
complements that discussion with the additional time-domain panels for
\(t_{\mathrm{wait}}=\SI{20}{\femto\second}\) and
\(t_{\mathrm{wait}}=\SI{40}{\femto\second}\), together with the supplementary
nonrephasing and absorptive frequency-domain figures.

\autoref{fig:rslts_monomer_2d_time_hom_inhom_t020_040} shows the monomer 2D
time-domain traces at
\(t_{\mathrm{wait}}=\SI{20}{\femto\second}\) and
\(\SI{40}{\femto\second}\), supplementing the \(t_{\mathrm{wait}}=0\) main-text
reference in~\autoref{fig:rslts_monomer_2d_time_hom_inhom}. Panels~(a) and~(b)
correspond to \(t_{\mathrm{wait}}=\SI{20}{\femto\second}\), and panels~(c)
and~(d) to \(t_{\mathrm{wait}}=\SI{40}{\femto\second}\). In each waiting-time
pair, the left panel shows the homogeneous case and the right panel shows the
inhomogeneous case. The qualitative time-domain structure is similar to the
\(t_{\mathrm{wait}}=0\) case discussed in the main text.

\begin{figure}[htbp]
    \centering
    \safeplaceholdergraphic[1.22\linewidth]{figures/rslts_monomer_2d_time_hom_inhom_t020_040}
    \caption{Monomer 2D time-domain traces for larger waiting times, with
    shared settings listed in~\autoref{subsec:results_monomer_model},
    \(t_{\mathrm{coh}},t_{\mathrm{det}}\in[0,\SI{600}{\femto\second}]\) and
    \(\Delta t=\SI{2}{\femto\second}\). Panels~(a) and~(b):
    \(t_{\mathrm{wait}}=\SI{20}{\femto\second}\). Panels~(c) and~(d):
    \(t_{\mathrm{wait}}=\SI{40}{\femto\second}\). In each waiting-time pair,
    the left panel shows the homogeneous case, \(\Delta_{\mathrm{inh}}=0\), and
    the right panel the inhomogeneous case,
    \(\Delta_{\mathrm{inh}}=\SI{200}{\per\centi\metre}\).}
    \label{fig:rslts_monomer_2d_time_hom_inhom_t020_040}
\end{figure}

\cref{fig:app_monomer_2d_freq_hom_inhom_nonrephasing,fig:app_monomer_2d_freq_hom_inhom_absorptive}
show the nonrephasing and absorptive frequency-domain waiting-time series. They
supplement the main-text rephasing series
in~\autoref{fig:rslts_monomer_2d_freq_hom_inhom_rephasing}. The separated
branches are included so that their diagonal and anti-diagonal widths can be
compared with the absorptive combination.

\begin{figure}[htbp]
    \centering
    \safeplaceholdergraphic[1.22\linewidth]{figures/rslts_monomer_2d_freq_hom_inhom_nonrephasing}
    \caption{Monomer 2D frequency-domain waiting-time series, nonrephasing
    component. Shared settings are listed
    in~\autoref{subsec:results_monomer_model}. Top, middle, and bottom rows correspond to
    \(t_{\mathrm{wait}}=\SI{0}{\femto\second}\),
    \(\SI{20}{\femto\second}\), and \(\SI{40}{\femto\second}\),
    respectively. (a): homogeneous case,
    \(\Delta_{\mathrm{inh}}=0\). (b): inhomogeneous case,
    \(\Delta_{\mathrm{inh}}=\SI{200}{\per\centi\metre}\).}
    \label{fig:app_monomer_2d_freq_hom_inhom_nonrephasing}
\end{figure}

\begin{figure}[htbp]
    \centering
    \safeplaceholdergraphic[0.75\linewidth]{figures/rslts_monomer_2d_freq_hom_inhom_absorptive}
    \caption{Monomer 2D frequency-domain waiting-time series, absorptive
    component. Shared settings are listed
    in~\autoref{subsec:results_monomer_model}. Top, middle, and bottom rows correspond to
    \(t_{\mathrm{wait}}=\SI{0}{\femto\second}\),
    \(\SI{20}{\femto\second}\), and \(\SI{40}{\femto\second}\),
    respectively. (a): homogeneous case,
    \(\Delta_{\mathrm{inh}}=0\). (b): inhomogeneous case,
    \(\Delta_{\mathrm{inh}}=\SI{200}{\per\centi\metre}\). Colour bars are
    omitted; the plotting convention normalises the range to \(1\).}
    \label{fig:app_monomer_2d_freq_hom_inhom_absorptive}
\end{figure}

\paragraph{Supplementary linewidths.}
The rephasing linewidths are discussed in the main text together with the
rephasing figure. The table here gives the nonrephasing and absorptive values
for the supplementary panels. In the inhomogeneous ensemble, the nonrephasing
branch has larger FWHM values than the rephasing branch. The absorptive FWHM
values are smaller than the separated branch magnitudes.

\begin{table}[htbp]
    \centering
    \scriptsize
    \caption{Diagonal and anti-diagonal FWHM values for the nonrephasing and
    absorptive monomer 2D frequency-domain spectra. Values are in
    \si{\per\centi\metre}; the rephasing values are listed
    in~\autoref{tab:results_monomer_2d_rephasing_linewidths}.}
    \label{tab:app_monomer_2d_branch_linewidths}
    \begin{tabular}{@{}llcrr@{}}
        \toprule
        Case & Branch & $t_{\mathrm{wait}}$ / fs & Diagonal & Anti-diagonal \\
        \midrule
        Homogeneous & Nonrephasing & 0  & 175.20 & 175.05 \\
        Homogeneous & Nonrephasing & 20 & 175.06 & 175.06 \\
        Homogeneous & Nonrephasing & 40 & 175.04 & 175.03 \\
        Homogeneous & Absorptive & 0  & 126.41 & 99.47 \\
        Homogeneous & Absorptive & 20 & 112.38 & 111.92 \\
        Homogeneous & Absorptive & 40 & 112.15 & 112.15 \\
        Inhomogeneous & Nonrephasing & 0  & 605.38 & 514.86 \\
        Inhomogeneous & Nonrephasing & 20 & 621.91 & 534.52 \\
        Inhomogeneous & Nonrephasing & 40 & 622.48 & 535.26 \\
        Inhomogeneous & Absorptive & 0  & 373.69 & 139.59 \\
        Inhomogeneous & Absorptive & 20 & 367.20 & 158.08 \\
        Inhomogeneous & Absorptive & 40 & 367.27 & 158.35 \\
        \bottomrule
    \end{tabular}
\end{table}

\subsection{Phenomenological and Microscopic Redfield Comparison}

\autoref{fig:app_monomer_legacy_fig4_t0_time} and
\autoref{fig:app_monomer_legacy_fig4_t0_freq} compare 2D inhomogeneous
monomer plots for \(N_{\mathrm{inhom}}=10\) and \(100\). Panel~(b) is visibly
smoother, while panel~(a) shows stronger jagged fluctuations from finite
sampling. These supplementary figures also test larger pulse amplitudes than in
the main text. Although these amplitudes populate excited states above \(1\%\)
(see~\autoref{app:sec:one_pulse_subtraction}), this does not affect the plotted
signals because the single-pulse polarisations are subtracted. Nevertheless,
the main-text discussion keeps the pulse amplitudes at \(0.2\) for consistency
with the paper.


\begin{figure}[htbp]
    \centering
    \safeplaceholdergraphic[1.22\linewidth]{figures/app_monomer_legacy_fig4_t0_time}
    \caption{Inhomogeneous phenomenological monomer 2D time-domain signals at
    \(t_{\mathrm{wait}}=\SI{0}{\femto\second}\). Shared settings are listed
    in~\autoref{subsec:results_monomer_model}. (a):
    \(N_{\mathrm{inhom}}=10\). (b): \(N_{\mathrm{inhom}}=100\). Both panels
    use pulse amplitudes \((0.005,0.005,0.005)\).}
    \label{fig:app_monomer_legacy_fig4_t0_time}
\end{figure}

\begin{figure}[htbp]
    \centering
    \safeplaceholdergraphic[0.95\linewidth]{figures/app_monomer_legacy_fig4_t0_freq}
    \caption{Inhomogeneous phenomenological monomer 2D frequency-domain spectra for
    time domain signals in~\autoref{fig:app_monomer_legacy_fig4_t0_time}. Shared
    settings are listed in~\autoref{subsec:results_monomer_model}. (a):
    \(N_{\mathrm{inhom}}=10\). (b): \(N_{\mathrm{inhom}}=100\).}
    \label{fig:app_monomer_legacy_fig4_t0_freq}
\end{figure}


\autoref{fig:app_monomer_legacy_redfield} shows a monomer signal from a microscopic Redfield simulation in the
laboratory frame, without the rotating-wave approximation. It also gives a
high-temperature bath example with a thermal initial state. The Ohmic bath
parameters are \(s=1\), \(T=\SI{11511}{\kelvin}\),
\(\omega_{\mathrm{c}}=100\bar{\omega}_0\), and
\(\alpha=0.001\bar{\omega}_0\). This parameter set gives
the same dephasing and population-decay rates as the phenomenological monomer
cases. Relative to the main-text monomer figures, parameters are
\(N_{\mathrm{inhom}}=50\) instead of \(100\) and \(t_{\mathrm{coh}},t_{\mathrm{det}}\in
[0,\SI{150}{\femto\second}]\) instead of
\([0,\SI{600}{\femto\second}]\).


\begin{figure}[htbp]
    \centering
    \safeplaceholdergraphic[1.22\linewidth]{figures/app_monomer_legacy_redfield}
    \caption{Two-dimensional signals for the inhomogeneous microscopic Redfield
    monomer at \(t_{\mathrm{wait}}=\SI{0}{\femto\second}\): time-domain signal
    (a) and frequency-domain signal (b).
    Shared settings are listed in~\autoref{subsec:results_monomer_model}. The
    plotted signal components are labelled in the figure. Parameters are
    \(N_{\mathrm{inhom}}=50\), a thermal initial state, bath temperature
    \SI{11511}{\kelvin}, and
    \(t_{\mathrm{coh}}=t_{\mathrm{det}}=\SI{150}{\femto\second}\).}
    \label{fig:app_monomer_legacy_redfield}
\end{figure}

\clearpage
\section{Dimer Figures}
\label{app:sec:paper_eqs_dimer}

This section gathers the dimer control figures and waiting-time series used in
the thesis discussion. The inhomogeneous waiting-time panels use
$N_{\mathrm{inhom}}=1000$. The \(t_{\mathrm{wait}}=0\) inhomogeneous control
figures in the main text use \(N_{\mathrm{inhom}}=10000\), as stated
in~\autoref{subsec:results_dimer_model}. The homogeneous companion sequence uses
\(\Delta_{\mathrm{inh}}=0\), so no disorder averaging is required.

\subsection{Waiting-Time Figures}

\crefrange{fig:app_dimer_fig4_time_t0_t16}{fig:app_dimer_fig4_time_t140_t310}
collect the dimer time-domain waiting-time panels that accompany the
frequency-domain waiting-time sequence in the main chapter.

\begin{figure}[htbp]
    \centering
    \safeplaceholdergraphic[1.22\linewidth]{figures/app_dimer_fig4_time_t0_t16}
    \caption{Dimer time-domain waiting-time panels for
    \(t_{\mathrm{wait}}=\SI{0}{\femto\second}\) (a) and
    \(t_{\mathrm{wait}}=\SI{16}{\femto\second}\) (b). Shared dimer settings
    are listed in~\autoref{subsec:results_dimer_model}.}
    \label{fig:app_dimer_fig4_time_t0_t16}
\end{figure}

\begin{figure}[htbp]
    \centering
    \safeplaceholdergraphic[1.22\linewidth]{figures/app_dimer_fig4_time_t30_t46}
    \caption{Dimer time-domain waiting-time panels for
    \(t_{\mathrm{wait}}=\SI{30}{\femto\second}\) (a) and
    \(t_{\mathrm{wait}}=\SI{46}{\femto\second}\) (b). Shared dimer settings
    are listed in~\autoref{subsec:results_dimer_model}.}
    \label{fig:app_dimer_fig4_time_t30_t46}
\end{figure}

\begin{figure}[htbp]
    \centering
    \safeplaceholdergraphic[1.22\linewidth]{figures/app_dimer_fig4_time_t62_t108}
    \caption{Dimer time-domain waiting-time panels for
    \(t_{\mathrm{wait}}=\SI{62}{\femto\second}\) (a) and
    \(t_{\mathrm{wait}}=\SI{108}{\femto\second}\) (b). Shared dimer settings
    are listed in~\autoref{subsec:results_dimer_model}.}
    \label{fig:app_dimer_fig4_time_t62_t108}
\end{figure}

\begin{figure}[htbp]
    \centering
    \safeplaceholdergraphic[1.22\linewidth]{figures/app_dimer_fig4_time_t140_t310}
    \caption{Dimer time-domain waiting-time panels for
    \(t_{\mathrm{wait}}=\SI{140}{\femto\second}\) (a) and
    \(t_{\mathrm{wait}}=\SI{310}{\femto\second}\) (b). Shared dimer settings
    are listed in~\autoref{subsec:results_dimer_model}.}
    \label{fig:app_dimer_fig4_time_t140_t310}
\end{figure}

\begin{figure}[htbp]
    \centering
    \safeplaceholdergraphic[0.98\linewidth]{figures/rslts_dimer_fig4_paper_eqs_waiting_time_absorptive}
    \caption{Waiting-time sequence for the coupled dimer, absorptive component.
    The shared dimer settings are listed in~\autoref{subsec:results_dimer_model}.
    Waiting times are the same as
    in~\autoref{fig:rslts_dimer_fig4_waiting_time_rephasing}. Colour bars are
    omitted; the plotting convention normalises the range to \(1\).}
    \label{fig:rslts_dimer_fig4_waiting_time_absorptive}
\end{figure}

\subsection{Homogeneous Companion Panels}
\autoref{fig:app_dimer_fig6_time} gives the homogeneous time-domain companion
panels. \autoref{fig:app_dimer_fig6_freq_absorptive} gives the corresponding
absorptive frequency-domain spectra.

\begin{figure}[htbp]
    \centering
    \safeplaceholdergraphic[0.70\linewidth]{figures/rslts_dimer_fig6_paper_eqs_homogeneous_time_all_components}
    \caption{Homogeneous dimer time-domain signals. Shared dimer settings are
    listed in~\autoref{subsec:results_dimer_model}. Top:
    \(t_{\mathrm{wait}}=\SI{46}{\femto\second}\). Bottom:
    \(t_{\mathrm{wait}}=\SI{62}{\femto\second}\). The panels show the rephasing,
    nonrephasing, and absorptive components.}
    \label{fig:app_dimer_fig6_time}
\end{figure}

\begin{figure}[htbp]
    \centering
    \safeplaceholdergraphic[0.74\linewidth]{figures/rslts_dimer_fig6_paper_eqs_homogeneous_freq_absorptive}
    \caption{Homogeneous dimer absorptive frequency-domain spectra
    for the time domain signals in~\autoref{fig:app_dimer_fig6_time}. Shared dimer
    settings are listed in~\autoref{subsec:results_dimer_model}. Colour bars
    are omitted; the plotting convention normalises the range to \(1\).}
    \label{fig:app_dimer_fig6_freq_absorptive}
\end{figure}

\subsection{Dimer Linewidths}

The main text uses the rephasing linewidths to compare compact homogeneous
peaks with diagonally broadened inhomogeneous peaks. The complete extracted
values are collected here. Each value is an FWHM in \si{\per\centi\metre}. The
widths were measured from local diagonal and anti-diagonal cuts through the two
diagonal dimer features of
\(|\tilde{E}(\omega_{\mathrm{coh}},\omega_{\mathrm{det}})|\).

\begin{table}[htbp]
    \centering
    \scriptsize
    \caption{Rephasing linewidths for the two diagonal dimer features. Values
    are FWHM in \si{\per\centi\metre}; coupling values in the figure/setup
    column are also in \si{\per\centi\metre}.}
    \label{tab:app_dimer_2d_rephasing_linewidths}
    \begin{tabular}{@{}lcccc@{}}
        \toprule
        Figure/setup & $t_{\mathrm{wait}}$ / fs & Peak & Diagonal & Anti-diagonal \\
        \midrule
        Fig.~\ref{fig:rslts_dimer_2d_freq_un_coupled} ($J=0$)   & 0  & 11 & 302.00 & 55.43 \\
        Fig.~\ref{fig:rslts_dimer_2d_freq_un_coupled} ($J=0$)   & 0  & 22 & 295.33 & 55.69 \\
        Fig.~\ref{fig:rslts_dimer_2d_freq_un_coupled} ($J=300$) & 0  & 11 & 241.64 & 54.16 \\
        Fig.~\ref{fig:rslts_dimer_2d_freq_un_coupled} ($J=300$) & 0  & 22 & 261.99 & 75.99 \\
        Fig.~\ref{fig:rslts_dimer_fig4_waiting_time_rephasing} ($J=500$) & 46 & 11 & 223.52 & 53.96 \\
        Fig.~\ref{fig:rslts_dimer_fig4_waiting_time_rephasing} ($J=500$) & 46 & 22 & 250.72 & 90.27 \\
        Fig.~\ref{fig:rslts_dimer_fig4_waiting_time_rephasing} ($J=500$) & 62 & 11 & 223.76 & 54.04 \\
        Fig.~\ref{fig:rslts_dimer_fig4_waiting_time_rephasing} ($J=500$) & 62 & 22 & 249.41 & 97.30 \\
        Fig.~\ref{fig:rslts_dimer_fig6_homogeneous_companion} ($J=500$)  & 46 & 11 & 70.27  & 70.23 \\
        Fig.~\ref{fig:rslts_dimer_fig6_homogeneous_companion} ($J=500$)  & 46 & 22 & 82.66  & 84.64 \\
        Fig.~\ref{fig:rslts_dimer_fig6_homogeneous_companion} ($J=500$)  & 62 & 11 & 70.00  & 70.01 \\
        Fig.~\ref{fig:rslts_dimer_fig6_homogeneous_companion} ($J=500$)  & 62 & 22 & 84.42  & 84.52 \\
        \bottomrule
    \end{tabular}
\end{table}

\autoref{tab:app_dimer_2d_rephasing_linewidths} shows the same trend for both
diagonal resonances. The homogeneous companion is compact in both directions.
With static disorder, the diagonal widths are much larger than the
anti-diagonal widths.

\autoref{tab:app_dimer_2d_nonreph_abs_linewidths} gives the corresponding
nonrephasing and absorptive linewidths.

\begin{table}[htbp]
    \centering
    \scriptsize
    \caption{Nonrephasing and absorptive linewidths for the two diagonal dimer
    features. Values are FWHM in \si{\per\centi\metre}; coupling values in the
    figure/setup column are also in \si{\per\centi\metre}.}
    \label{tab:app_dimer_2d_nonreph_abs_linewidths}
    \begin{tabular}{@{}llcccc@{}}
        \toprule
        Figure/setup & Branch & $t_{\mathrm{wait}}$ / fs & Peak & Diagonal & Anti-diagonal \\
        \midrule
        \autoref{fig:rslts_dimer_2d_freq_un_coupled}, $J=0$ & Nonrephasing & 0 & $11$ & 382.08 & 281.18 \\
        \autoref{fig:rslts_dimer_2d_freq_un_coupled}, $J=0$ & Nonrephasing & 0 & $22$ & 365.81 & 318.67 \\
        \autoref{fig:rslts_dimer_2d_freq_un_coupled}, $J=0$ & Absorptive & 0 & $11$ & 295.68 & 30.72 \\
        \autoref{fig:rslts_dimer_2d_freq_un_coupled}, $J=0$ & Absorptive & 0 & $22$ & 286.33 & 30.80 \\
        \autoref{fig:rslts_dimer_2d_freq_un_coupled}, $J=300$ & Nonrephasing & 0 & $11$ & 352.02 & 268.22 \\
        \autoref{fig:rslts_dimer_2d_freq_un_coupled}, $J=300$ & Nonrephasing & 0 & $22$ & 388.31 & 301.26 \\
        \autoref{fig:rslts_dimer_2d_freq_un_coupled}, $J=300$ & Absorptive & 0 & $11$ & 221.65 & 29.67 \\
        \autoref{fig:rslts_dimer_2d_freq_un_coupled}, $J=300$ & Absorptive & 0 & $22$ & 240.64 & 45.74 \\
        \autoref{fig:rslts_dimer_fig4_waiting_time_rephasing}, $J=500$ & Nonrephasing & 46 & $11$ & 201.09 & 163.41 \\
        \autoref{fig:rslts_dimer_fig4_waiting_time_rephasing}, $J=500$ & Nonrephasing & 46 & $22$ & 200.84 & 188.22 \\
        \autoref{fig:rslts_dimer_fig4_waiting_time_rephasing}, $J=500$ & Absorptive & 46 & $11$ & 217.54 & 27.45 \\
        \autoref{fig:rslts_dimer_fig4_waiting_time_rephasing}, $J=500$ & Absorptive & 46 & $22$ & 241.16 & 47.77 \\
        \autoref{fig:rslts_dimer_fig4_waiting_time_rephasing}, $J=500$ & Nonrephasing & 62 & $11$ & 283.60 & 249.54 \\
        \autoref{fig:rslts_dimer_fig4_waiting_time_rephasing}, $J=500$ & Nonrephasing & 62 & $22$ & 249.77 & 261.01 \\
        \autoref{fig:rslts_dimer_fig4_waiting_time_rephasing}, $J=500$ & Absorptive & 62 & $11$ & 210.98 & 29.21 \\
        \autoref{fig:rslts_dimer_fig4_waiting_time_rephasing}, $J=500$ & Absorptive & 62 & $22$ & 202.95 & 63.63 \\
        \autoref{fig:rslts_dimer_fig6_homogeneous_companion}, $J=500$ & Nonrephasing & 46 & $11$ & 71.93 & 71.46 \\
        \autoref{fig:rslts_dimer_fig6_homogeneous_companion}, $J=500$ & Nonrephasing & 46 & $22$ & 79.03 & 82.28 \\
        \autoref{fig:rslts_dimer_fig6_homogeneous_companion}, $J=500$ & Absorptive & 46 & $11$ & 49.53 & 30.02 \\
        \autoref{fig:rslts_dimer_fig6_homogeneous_companion}, $J=500$ & Absorptive & 46 & $22$ & 98.68 & 44.63 \\
        \autoref{fig:rslts_dimer_fig6_homogeneous_companion}, $J=500$ & Nonrephasing & 62 & $11$ & 70.12 & 70.13 \\
        \autoref{fig:rslts_dimer_fig6_homogeneous_companion}, $J=500$ & Nonrephasing & 62 & $22$ & 77.06 & 76.62 \\
        \autoref{fig:rslts_dimer_fig6_homogeneous_companion}, $J=500$ & Absorptive & 62 & $11$ & 33.76 & 42.13 \\
        \autoref{fig:rslts_dimer_fig6_homogeneous_companion}, $J=500$ & Absorptive & 62 & $22$ & 39.88 & 60.40 \\
        \bottomrule
    \end{tabular}
\end{table}

The nonrephasing branch has larger FWHM values in the inhomogeneous cases,
especially along the anti-diagonal direction. The absorptive entries should be
read cautiously because they describe the combined branch rather than either
separated branch.
